% -----------------------------------------------------------------------------------------------------
\section{算符的 BRST 分类：非定域缺口}\label{sec:front-brst}

\begin{remark}[为什么需要 BRST：K 胶子与非定域缺口]\label{rem:front-k-gluon}
在 Feynman 规范中，DIS 的领头区域含任意多个附着于硬部的纵向极化（“K”）胶子。图形化的非阿贝尔 Ward 恒等式证明\cite{Labastida:1984gy}（\cite[第~11.9~节]{Collins2011}）只在严格共线极限下成立，余项在领头幂次不为零：纯 K 胶子项经由胶子横动量对胶子密度有贡献（显式计算见\cite{Collins:2008sg}），这是构造规范不变胶子 PDF 所需的非阿贝尔效应。Collins 因此改走 BRST 路线\cite[第~11.3.2~节，第~407--410~页]{Collins2011}：对离壳 Green 函数 $T^{\mu\nu}_{\rm off\text{-}shell}(q)=\frac1{4\pi}\int\dd^4z\,\ee^{\ii q\cdot z}\langle0|T\,\text{fields}\,j^\mu(z/2)j^\nu(-z/2)|0\rangle$（其式 (11.11)），领头区域给出 $\sum_iC_i\otimes\langle0|T\,\text{fields}\,O_i|0\rangle$（式 (11.12)），其中 $O_i$ 先验地是光线上“arbitrarily complicated multilocal operators”；Green 函数（式 (11.13)）与流的 BRST 不变性给出 $\delta_{\rm BRST}\sum_iC_i\otimes O_i=0$（式 (11.17)）。对\emph{定域}算符，BRST 不变 $=$ 规范不变 $\oplus$ (A) BRST 变分 $\oplus$ (B) 由运动方程为零的算符（\cite{Joglekar:1975nu}，以及 Collins 所引的\cite{Joglekar:1976eb,Joglekar:1976pe,Nakanishi:1990qm}；作为定域 BRST 上同调的定理见\cite{Henneaux:1993jn,Barnich:2000zw}；Ward 恒等式与背景场规范见\cite{Kluberg-Stern:1974iel,Kluberg-Stern:1974nmx,Kluberg-Stern:1975ebk}，最早的显式规范变异抵消项见\cite{DixonTaylor1974}）——核心陈述是\cref{prop:appJ-local-operators}（[R, $R_{\mathrm{PT}}$]）；A、B 两类在物理矩阵元中为零。缺口的原文（第~410~页）：“Unfortunately, the published proofs that BRST-invariant operators are either gauge invariant or are in one of classes A or B apply as written to local operators. It is natural that the result also applies to the non-local operators we use in factorization. But, as far as I know, no proof has been given. For the purposes of the discussion, I will assume the result is true, and leave the proof (or refutation) to future research.”超出单圈的离壳计算需要 A、B 类算符\cite{Hamberg:1991qt,Collins:1994ee}；规范变异抵消项的全圈结构见\cite{Falcioni:2022fdm}（显式基到四圈）与\cite{Gehrmann:2023ksf}（三圈单态反常维度）（均据摘要）。光锥规范 $A^+=0$ 消去 K 胶子，但 Collins 写道，“given the known problems with $A^+=0$ gauge, it would be nice to have a proof that does not rely on the gauge”。即使承认规范不变性，Wilson 线的路径也要由其余要求确定：共线 PDF 是直的类光线，TMD 则是依赖过程的订书钉形路径。
\end{remark}

\begin{proposition}[双局域 PDF 算符 BRST 分类的矩约化 \tagA]\label{prop:front-brst-bilocal}\hfill\break
令 $O_q(w^-)=\bar\psi(w)\gamma^+W[w,0]\psi(0)$，$w=(0^+,w^-,\mathbf 0_T)$，$O_g(w^-)$ 为相应的胶子算符（\cref{def:pdf-lightray}）。在物理（在壳、BRST 不变）矩阵元中，于重整化微扰论的每一阶，$O^R=Z_{OO}\otimes O_{(0)}$，且
\begin{equation}\label{eq:front-moment-reduction}
\int_0^1\dd z\,z^{N-1}Z_{OO}(z)=Z^{(N)}_{\rm twist\text{-}2}\qquad\text{对一切 }N ,
\end{equation}
BRST 恰当类（A）与运动方程类（B）的混入不出现，$Z_{OO}$ 由其矩唯一确定（紧支；对 $C$ 偶、$C$ 奇组合 $q\pm\bar q$ 与胶子只出现偶数或奇数的 $N$，唯一性由\cref{thm:appK-muntz}给出，至多相差 $z=0$ 处的分布）。因此三角混合（\cite[式~(11.20)]{Collins2011}）
\begin{equation}\label{eq:front-triangular}
\begin{pmatrix}O\\A\\B\end{pmatrix}=\begin{pmatrix}Z_{OO}&Z_{OA}&Z_{OB}\\0&Z_{AA}&Z_{AB}\\0&0&Z_{BB}\end{pmatrix}\begin{pmatrix}O_{(0)}\\A_{(0)}\\B_{(0)}\end{pmatrix}
\end{equation}
限制在物理矩阵元上对双局域 PDF 算符成立。所用输入：$\bar\psi\gamma^+(\ii D^+)^{N-1}\psi$ 与 $F^{+\mu}(\ii D^+)^{N-2}F^+{}_\mu$ 的定域分类（指标全为“$+$”，因而自动对称无迹，恰为扭度 2 算符，\cref{prop:appJ-local-operators}）；前向矩阵元逐阶支撑在 $[-1,1]$（\cref{thm:pdf-support-pt}）；裸非定域算符是裸定域算符的母函数（\cref{prop:appL-bare}；对紧支分布精确，因为其 Fourier 变换是整函数，Taylor 系数即矩）；重整化的矩关系（\cref{prop:appL-renormalized}）。（本命题只涉及算符与部分子态。）
% VF-FRO-15
\uses{def:pdf-lightray, prop:appJ-local-operators, thm:pdf-support-pt, prop:appL-bare, prop:appL-renormalized, thm:appK-muntz}
\end{proposition}

\begin{proofsketch}
(1)~裸双局域算符的矩就是裸定域扭度 2 算符（\cref{prop:appL-bare}）。(2)~由定域分类，$[O_N]=Z_NO_N^{\rm bare}+(\text{A 类})+(\text{B 类})$，且混合是三角的；在物理矩阵元中只剩 $Z_N\equiv Z^{(N)}_{\rm twist\text{-}2}$（\cref{prop:appJ-local-operators}）。(3)~紧支意味着双局域算符的物理矩阵元由其全部矩决定（\cref{thm:pdf-support-pt}）。(4)~在每个 $C$ 宇称扇区（单态为夸克—胶子矩阵）逐阶存在 $[0,1]$ 上的分布 $Z_{OO}$，其 Mellin 矩为 $Z_N$（单圈由 $[(1+z^2)/(1-z)]_+$ 构成）；唯一性：扇区内由\cref{thm:appK-muntz}在 $(0,1]$ 上唯一，在 $[-1,1]$ 上取全部 $N$ 则只需紧支。(5)~重整化的矩关系（\cref{prop:appL-renormalized}）把 (1)--(4) 合成 $O^R=Z_{OO}\otimes O_{(0)}$，即\cref{eq:front-moment-reduction}与限制在物理矩阵元上的\cref{eq:front-triangular}。单圈情形可以逐图核对：由算符定义算出的单圈核的 Mellin 矩等于定域扭度 2 算符的单圈反常维度（第八章）。
\end{proofsketch}

\paragraph{地位与缺口。}[A]。这一约化有先行工作：一般非前向形式是\cite[式~(2.4)、(2.6)、(3.1)、(3.2)]{Braun:2019qtp}（光线算符作为定域算符的母函数，重整化的光线算符由重整化的定域算符同样定义，Joglekar--Lee 分解），思想可追溯到\cite{Balitsky:1987bk}；本书不把它当作新结果。所有输入在本书中都是 [R] 或数学，但我们没有逐行核查附录~L 中\cref{prop:appL-renormalized}的证明不预设非定域算符的分类（循环性问题），这是最弱的一步（代数约化，误差类型不适用）；因此不升级为 [R]。保留：对维数正规化中前向在壳部分子矩阵元，BRST 类与运动方程类不一定为零。它\emph{不}覆盖：因子化证明\emph{内部}的分类一步（规范不变性施加之前、场位于不同光锥位置的多局域算符，\cref{conj:front-brst-multilocal}）；TMD 算符（订书钉路径与快度发散，是另一个问题）；光线算符本身的离壳重整化（物理预言不需要）。

\begin{conjecture}[多局域光线算符的 BRST 分类 \tagC]\label{conj:front-brst-multilocal}
在离壳 Green 函数的领头区域分解（\cite[式~(11.12)]{Collins2011}）中出现的那类、支撑在一条类光线上、鬼数为零的 BRST 不变多局域算符，都等于一个规范不变算符、一个 BRST 变分（A 类）与一个由运动方程为零的算符（B 类）之和。
\uses{def:prelim-brst, prop:appJ-local-operators}
\end{conjecture}

\paragraph{地位与缺口。}[C]：Collins 假定了它（\cref{rem:front-k-gluon}中的原文）；据我们截至 2026 年 10 月的文献检索（INSPIRE 与网络检索），尚未见非定域（多局域）BRST 分类的证明。在因子化证明中它以被明确假定的一步\cref{asm:appN-brst-nonlocal}（[C]）出现。自然的工具是用辅助场局域化（把 Wilson 线当作一个场），如同类空 Wilson 线算符的可重整性论证\cite{Ji:2017oey,Ishikawa:2017faj,Green:2017xeu}（\cref{prop:lattice-renormalizability}，\cref{asm:lamet-aux-renorm}）。“Collins（2011）证明了非定域算符的 BRST 分类”是过度表述；附录~J 的\cref{rem:appJ-nonlocal}指向这里与\cref{prop:front-brst-bilocal}。
